\documentclass[12pt, a4paper]{ctexart}
\usepackage{geometry} 
\usepackage{amsmath}     
\usepackage{graphicx}    
\usepackage{booktabs}    
\usepackage{tabularx}    
\usepackage{xcolor}      
\usepackage{caption} 
\usepackage{fancyhdr}    
\usepackage{hyperref}
\usepackage{enumitem}
\usepackage{tikz}
\usepackage[most]{tcolorbox}

% --- 自定义 tcolorbox 样式 ---
% 经典引用框（红色主题）
\newtcolorbox{quotebox}[1][]{
  colback=red!3, colframe=red!50!black,
  fonttitle=\bfseries, title=#1,
  arc=2mm, boxrule=0.8pt,
  left=8pt, right=8pt, top=6pt, bottom=6pt,
  breakable
}
% 蓝色引用框
\newtcolorbox{bluequotebox}[1][]{
  colback=blue!3, colframe=blue!50!black,
  fonttitle=\bfseries, title=#1,
  arc=2mm, boxrule=0.8pt,
  left=8pt, right=8pt, top=6pt, bottom=6pt,
  breakable
}
% 章节小结框（灰色主题）
\newtcolorbox{summarybox}[1][本节小结]{
  colback=gray!8, colframe=gray!60,
  fonttitle=\bfseries\kaishu, title=#1,
  arc=1.5mm, boxrule=0.6pt,
  left=8pt, right=8pt, top=6pt, bottom=6pt
}

% --- 定义上标引用命令 ---
\newcommand{\upcite}[1]{\textsuperscript{\cite{#1}}}

\geometry{left=2.5cm, right=2.5cm, top=3cm, bottom=3cm}
\hypersetup{
    colorlinks=true,
    linkcolor=blue!70!black,
    citecolor=blue!70!black,
    urlcolor=blue!70!black
}
\captionsetup{font=small, labelfont=bf}
\graphicspath{{figure/}}

% --- 图表编号带章节号 ---
\numberwithin{table}{section}
\numberwithin{figure}{section}
\renewcommand{\thetable}{\thesection-\arabic{table}}
\renewcommand{\thefigure}{\thesection-\arabic{figure}}

% --- 章节标题格式 ---
\ctexset{
  section = {
    format = \heiti\zihao{-3}\centering,
    beforeskip = 1.5ex plus 0.5ex,
    afterskip = 1ex plus 0.3ex,
    name = {,},
    number = \chinese{section},
    aftername = \quad
  },
  subsection = {
    format = \heiti\zihao{4}\raggedright,
    beforeskip = 1ex plus 0.3ex,
    afterskip = 0.8ex plus 0.2ex,
    name = {（,）},
    number = \chinese{subsection},
    aftername = \quad
  }
}

% --- 页眉页脚设置 ---
\pagestyle{fancy}
\fancyhf{} % 清除所有默认设置
\renewcommand{\headrulewidth}{0.5pt}

% 居中页眉：作者等：稿件题目
\fancyhead[C]{\zihao{5} 杨林涛：从”走苏联的路”到”走自己的路”}
% 页脚居中页码
\fancyfoot[C]{\zihao{-5} —~\thepage~—}
% 设置全局1.5倍行距
\linespread{1.5}

% --- 定义前页样式（目录等用罗马数字） ---
\fancypagestyle{frontmatter}{
  \fancyhf{}
  \renewcommand{\headrulewidth}{0.5pt}
  \fancyhead[C]{\zihao{5} 杨林涛：从”走苏联的路”到”走自己的路”}
  \fancyfoot[C]{\zihao{-5} —~\thepage~—}
}

% --- 定义摘要框样式 ---
\newtcolorbox{abstractbox}[1][]{
  colback=gray!4, colframe=gray!40,
  arc=1.5mm, boxrule=0.4pt,
  left=10pt, right=10pt, top=8pt, bottom=8pt,
  breakable, #1
}

\begin{document}
\thispagestyle{empty} % 封面页不显示页码

\begin{center}
    \vspace*{1em}
    \begin{tcolorbox}[
      colback=red!2, colframe=red!40!black,
      arc=3mm, boxrule=1pt,
      width=0.92\textwidth,
      top=15pt, bottom=15pt,
      halign=center
    ]
    {\heiti \zihao{-3} \textbf{从”走苏联的路”到”走自己的路”}} \\[0.8em]
    {\kaishu \zihao{-4} ——探讨中国社会主义建设道路的自主性探索历程}
    \end{tcolorbox}
    \vspace{2em}

    {\zihao{-4}
    \renewcommand{\arraystretch}{1.5}
    \begin{tabular}{rl}
        \textbf{课程名称：} & 社会主义发展史 \\
        \textbf{所属院校：} & 浙江大学 \\
        \textbf{学生姓名：} & 杨林涛 \\
        \textbf{学生学号：} & 3250106024 \\
        \textbf{完成日期：} & 2026年3月2日（当前日期） \\
    \end{tabular}
    }
\end{center}
\vspace{2em}
{\color{red!40!black}\hrule height 0.8pt}
\vspace{2em}
\begin{abstractbox}
    {\kaishu \zihao{-4}
    \centerline{\heiti\textbf{【摘\quad 要】}}
    \vspace{0.3em}
    本文整理中国建设社会主义的过程，最初学习苏联，后来建立中国特色社会主义。文章分析这个过程的原因，国家刚建立时，西方国家封锁中国，经济非常落后，中国借鉴苏联的方法，借此建立工业基础。后来，苏联方法出现问题，重工业发展不平衡，权力过于集中。毛泽东等领导人发表《论十大关系》，中共召开第八次全国代表大会，中国开始自己探索建设道路。中国在二十世纪六七十年代遇到很大困难，但是，当时的经验和思考帮助了后来的改革。中共十一届三中全会召开，中国放弃单一的计划经济，建立社会主义市场经济，成功找到中国特色社会主义道路。文章指出，中国一直坚持自己做决定，这是这条道路的主要特点，也是新时代推进中国式现代化的重要动力和基本原则。
    \vspace{0.5em}
    \newline
    \textbf{【关键词】} 社会主义建设；苏联模式；自主探索；改革开放；中国特色社会主义
    }
\end{abstractbox}

\vspace{1.5em}

\begin{abstractbox}
    {\zihao{-4}
    \centerline{\textbf{From ``Following the Soviet Path'' to ``Forging Our Own Path''}}
    \vspace{0.2em}
    \centerline{\zihao{5}\textit{--- Exploring the History of China's Autonomous Quest for a Socialist Construction Road}}
    \vspace{0.5em}
    \centerline{\textbf{[Abstract]}}
    \vspace{0.3em}
    {\zihao{5} This paper examines China's transition from emulating the Soviet Union to establishing socialism with Chinese characteristics. In its early years, the newly founded People's Republic faced Western blockades and severe economic backwardness, compelling the adoption of the Soviet model to build an industrial foundation. However, the inherent deficiencies of this model---imbalanced emphasis on heavy industry and excessive centralization of power---soon became apparent. Mao Zedong's \textit{On the Ten Major Relationships} and the Eighth National Congress of the CPC marked the beginning of China's independent exploration. Despite significant setbacks during the 1960s and 1970s, the experiences and reflections of that period laid the groundwork for subsequent reforms. The Third Plenary Session of the Eleventh Central Committee initiated reform and opening-up, leading China to abandon rigid central planning in favor of a socialist market economy. This paper argues that the spirit of autonomous decision-making constitutes the defining feature of China's development path and serves as a fundamental driving force for advancing Chinese-style modernization in the new era.}
    \vspace{0.5em}
    \newline
    \textbf{[Keywords]} Socialist Construction; Soviet Model; Autonomous Exploration; Reform and Opening-up; Socialism with Chinese Characteristics
    }
\end{abstractbox}
\vspace{1em}
\newpage
\pagenumbering{roman}  % 目录页使用罗马数字
\pagestyle{frontmatter}
\tableofcontents

\newpage
\listoffigures
\vspace{2em}
\listoftables
\newpage

% --- 正文开始 ---
\pagenumbering{arabic}  % 正文使用阿拉伯数字页码
\pagestyle{fancy}
\songti \zihao{-4}

\section{绪论}
\subsection{研究背景与意义}
马克思等早期学者没有说明如何在经济落后的东方大国建设社会主义，一九四九年新中国成立，帝国主义国家封锁中国，国内经济非常破败。中国共产党为了保护新的政权和发展国家工业，决定学习苏联，这就是“走苏联的路”\upcite{ref1}。后来苏联方法的问题逐渐暴露，管理制度死板导致经济发展不平衡，中国共产党人认识到不能完全照抄外国做法，其不能解决中国复杂的实际问题。所以，文章研究中国从学习苏联到自己寻找道路的过程，这个研究帮助我们理解马克思主义如何结合中国实际，也帮助我们在新时代坚持发展中国式现代化。

\subsection{核心概念界定}
文章研究“苏联模式”和“自己做主”。“苏联模式”\footnote{“苏联模式”又称”斯大林模式”，学界对其内涵存在广义与狭义之分。狭义仅指斯大林时期形成的高度集中的政治经济体制；广义则涵盖苏联社会主义建设的整套理论、制度与实践模式。本文取其狭义。}特别指斯大林时期的做法，国家通过死板的命令管理经济，优先建设重工业，政治上权力过度集中在少数人手中\upcite{ref2}。“自己做主”指中国共产党坚持马克思主义，根据中国人口多、农业基础薄弱的真实情况做事，自己思考和决定，在实际工作中改变传统的做法，这是马克思主义结合中国实际的核心内容。

\subsection{研究思路与方法}
文章使用马克思主义的唯物主义方法研究问题，把中国建设社会主义的过程放在具体的历史时间里分析。详细分析中国早期学习苏联的过程，在困难中前进的过程，以及在新时期成功改变做法的过程。文章研究《论十大关系》、中共八大、十一届三中全会和邓小平南方讲话等重要事件，通过这些事件，说明中国共产党人如何自己决定国家发展方向。

% --- 时间线图：中国社会主义建设道路演进 ---
\vspace{1em}
\begin{figure}[htbp]
\centering
\begin{tikzpicture}[x=0.42cm, font=\zihao{6}]
  % 三个阶段背景色块
  \fill[red!10, rounded corners=2pt] (0,-0.6) rectangle (7,0.6);
  \fill[orange!10, rounded corners=2pt] (7,-0.6) rectangle (29,0.6);
  \fill[green!10, rounded corners=2pt] (29,-0.6) rectangle (43,0.6);

  % 主时间轴
  \draw[->, thick, gray!70] (-0.5,0) -- (44,0);

  % 阶段标注（轴下方）
  \node[red!60!black, font=\zihao{7}\bfseries] at (3.5,-1) {学习苏联};
  \node[orange!70!black, font=\zihao{7}\bfseries] at (18,-1) {曲折探索};
  \node[green!60!black, font=\zihao{7}\bfseries] at (36,-1) {改革开放};

  % 节点与标注（交替上下避免拥挤）
  \fill[red!70!black] (0,0) circle (3pt);
  \node[above, red!70!black, text width=1.5cm, align=center] at (0,0.4) {1949\\新中国\\成立};

  \fill[red!70!black] (4,0) circle (3pt);
  \node[below, red!70!black, text width=1.5cm, align=center] at (4,-1.5) {1953\\一五计划};

  \fill[orange!80!black] (7,0) circle (3pt);
  \node[above, orange!80!black, text width=2cm, align=center] at (7,0.4) {1956\\论十大\\关系/八大};

  \fill[orange!80!black] (9,0) circle (3pt);
  \node[below, orange!80!black, text width=1.5cm, align=center] at (9,-1.5) {1958\\大跃进};

  \fill[orange!80!black] (17,0) circle (3pt);
  \node[above, orange!80!black, text width=1.5cm, align=center] at (17,0.4) {1966\\文化大\\革命};

  \fill[green!60!black] (29,0) circle (3pt);
  \node[below, green!60!black, text width=2cm, align=center] at (29,-1.5) {1978\\改革开放};

  \fill[green!60!black] (36,0) circle (3pt);
  \node[above, green!60!black, text width=2cm, align=center] at (36,0.4) {1992\\社会主义\\市场经济};

  \fill[green!60!black] (43,0) circle (3pt);
  \node[below, green!60!black, text width=1.5cm, align=center] at (43,-1.5) {新时代\\中国式\\现代化};
\end{tikzpicture}
\caption{中国社会主义建设道路演进时间线}
\end{figure}
\vspace{1em}

\section{早期模仿与反思：建国初期”走苏联的路”的历史必然与局限}
\subsection{“一边倒”战略与苏联模式的引入}
新中国刚建立时，国内和国外的情况非常危险，在国外，美国带领西方资本主义国家孤立新中国，他们不和中国做生意，用军队包围中国。在国内，中国经历了很多年的战争，国家的经济非常落后，几乎没有现代工业。世界当时分为两个对立的阵营，毛泽东为了保护国家安全，也为了快速建设工业，决定全面依靠社会主义国家，全面学习苏联的建设经验，这是当时唯一的选择\upcite{ref3}。

在这种情况里，中国制定并执行了第一个五年计划，这个计划的时间是一九五三年到一九五七年，主要内容是建设一百五十六个大型工程\footnote{"156项工程"是中苏两国政府于1950年代商定的援华建设项目，涵盖钢铁、煤炭、电力、机械、化工、军工等领域。因中途调整，实际施工建设的项目为150项，但学界沿用"156项工程"这一通称。}，苏联帮助中国建设这些工程。中国决定优先建设重工业，逐步建立由国家完全控制的经济管理方式，这种做法在很短的时间里集中了全国的力量。中国用极短的时间建立了很多现代工业工厂，中国开始制造飞机，中国开始制造汽车，中国开始制造大型机器，中国由此建立了国家工业的基础\upcite{ref4}。

\subsection{苏联模式弊端在中国的显现}
但是，苏联的做法本身存在很多问题，中国全面开始建设社会主义，苏联做法不适合中国的情况，这些管理问题开始明显地暴露出来\upcite{ref5}。

第一个问题是国家经济不同部分发展很不平衡，苏联的做法是压低农民卖粮食的价格，同时提高农民买工业产品的价格，用这个方法强制收集建设重工业的钱。在这种做法下，国家要求农民把粮食全部卖给国家，并负责分配这些粮食，这保证了工业建设有足够的物资。但是，国家为了快速发展重工业，损害了农业和轻工业的发展，时间长了，农业生产的发展非常慢，这直接导致城市和农村老百姓的生活无法改善。

第二个问题是经济管理权力过于集中，苏联的做法是国家下达死板的命令，中央政府管理所有的具体事情，管理得太严格。这种做法严重降低了地方政府的工作热情，降低了工厂的工作热情，也降低了普通工人的工作热情，具体的工厂没有权力决定自己的工作。中国是一个面积很大的国家，各地的情况完全不同，作为一个农业国家，强行使用这种死板的命令做法，带来的坏处越来越明显。

\vspace{0.5em}
\begin{table}[htbp]
\centering
\caption{苏联模式在中国暴露的主要弊端}
\zihao{5}
\begin{tabularx}{0.92\textwidth}{>{\bfseries\centering\arraybackslash}p{2.2cm} X X}
\toprule
\textbf{问题领域} & \textbf{苏联做法} & \textbf{导致的后果} \\
\midrule
经济结构失衡 & 压低农产品价格，强制积累资金优先发展重工业 & 农业和轻工业发展缓慢，城乡人民生活无法改善 \\
\addlinespace
管理权力集中 & 中央下达死板命令，统管所有具体经济事务 & 地方、工厂、工人积极性严重下降，无法因地制宜 \\
\bottomrule
\end{tabularx}
\end{table}

\begin{summarybox}
苏联模式帮助中国在短期内建立了工业基础，但其经济结构失衡和管理体制僵化的弊端在中国日益凸显，迫使中国共产党人开始反思和寻找新的建设道路。
\end{summarybox}

\section{自主探索的起步：对中国特色社会主义建设规律的初步认识}
\subsection{《论十大关系》：探索本国建设道路的开篇之作}
一九五六年是新中国历史上的一个关键年份，中国在那一年基本完成了社会主义改造，作为一个位于东方的农业国家，如何全面建设社会主义？这是一个重大的历史问题。那一年的四月，毛泽东发表了《论十大关系》的讲话\footnote{《论十大关系》最初是毛泽东于1956年4月25日在中共中央政治局扩大会议上的讲话，后于同年5月2日在最高国务会议上再次阐述。该讲话全文于1976年12月26日在《人民日报》首次公开发表，收入《毛泽东选集》第五卷。}，这说明中国共产党开始寻找适合中国的建设方法\upcite{ref6}。

在改变经济结构方面，毛泽东在讲话中首先讨论重工业、轻工业和农业之间的关系，苏联的做法只看重重工业，不关心农业和轻工业，甚至强迫农民交出大部分粮食，严重损害农民的利益。毛泽东针对这些教训提出新的要求，他明确提出，中国必须发展农业和轻工业，用发展农业和轻工业的办法来帮助发展重工业。这说明中国深刻认识到苏联做法的错误，不再完全跟着苏联走，中国共产党自己决定如何安排农业、轻工业和重工业的发展速度。在调整政治关系方面，讲话提出我们要让党内和党外的人都发挥作用，我们要让国内和国外的人都发挥作用，我们要把中国建设成一个强大的国家。这个讲话改变了过去只看重阶级斗争的想法，让更多的人愿意一起建设社会主义\upcite{ref7}。

\subsection{中共八大：主要矛盾的精准论断与路线确立}
《论十大关系》初步改变了人们的想法，一九五六年九月，中共召开了第八次全国代表大会，这次大会把中国自己寻找道路的做法变成了正式的制度和计划。大会根据历史唯物主义的观点，准确说明了当时中国社会的主要问题，指出，人民需要经济和文化快速发展，但是，当时的经济和文化水平不能满足人民的需要，这就是国内的主要矛盾。这个结论要求党和国家改变工作重点，停止把阶级斗争放在第一位，开始专心保护和发展生产工作\upcite{ref8}。

在经济管理的具体做法上，中共八大也表现出自己创造新方法的精神，当时国家对经济活动管理得太多，管理得太死板。陈云在大会上提出了一个重要的想法——“三个主体，三个补充”：

\begin{bluequotebox}[陈云：”三个主体，三个补充”]
\begin{description}[style=nextline, leftmargin=1.5em, font=\bfseries]
  \item[经营形式] 国家和集体经营为主体，个人经营为补充
  \item[生产计划] 计划生产为主体，自由生产为补充
  \item[市场形式] 国家市场为主体，自由市场为补充
\end{description}
\end{bluequotebox}

\noindent 这个想法改变了完全控制经济的做法，开始允许市场发挥一点作用，这是中国共产党最早自己思考如何用多种方式管理社会主义经济。

\vspace{0.5em}
\begin{table}[htbp]
\centering
\caption{中国自主探索社会主义建设道路的关键里程碑}
\zihao{5}
\begin{tabularx}{0.92\textwidth}{>{\centering\arraybackslash}p{2cm} >{\centering\arraybackslash}p{3cm} X}
\toprule
\textbf{年份} & \textbf{事件} & \textbf{历史意义} \\
\midrule
1956年4月 & 《论十大关系》 & 开篇之作，首次提出不照搬苏联，自主安排经济结构 \\
\addlinespace
1956年9月 & 中共八大 & 准确判断主要矛盾，工作重心转向经济建设 \\
\addlinespace
1978年 & 真理标准大讨论 & 打破思想束缚，为改革开放奠定认识基础 \\
\addlinespace
1978年12月 & 十一届三中全会 & 停止阶级斗争路线，确立改革开放的历史起点 \\
\addlinespace
1982年 & 中共十二大 & 邓小平正式提出"走自己的路，建设有中国特色的社会主义" \\
\addlinespace
1992年初 & 邓小平南方谈话 & 突破计划与市场之争，明确市场是经济手段 \\
\addlinespace
1992年10月 & 中共十四大 & 正式确立社会主义市场经济体制目标 \\
\bottomrule
\end{tabularx}
\end{table}

\begin{summarybox}
《论十大关系》和中共八大标志着中国开始自主探索社会主义建设道路，提出了不同于苏联的发展思路，在经济结构调整和主要矛盾判断上迈出了关键一步。
\end{summarybox}

\section{曲折中的前进：自主探索中的失误与宝贵财富}
我们在破除对苏联做法盲目相信的过程中，经历了很多困难，付出了巨大的代价。我们坚持唯物主义要求，必须客观全面地看待这段时间的错误和成绩。

\subsection{急躁冒进与探索轨道的偏离}
二十世纪五十年代末到七十年代，全党看到国家落后，非常着急改变落后的情况，也缺少大规模建设国家的经验，在寻找道路时犯了严重的错误。

在指导经济建设的想法上，部分人仅凭自己的主观想法做事，认为只要有决心就能做成任何事。国家发动了“大跃进”运动，试图用发动大量老百姓快速工作的方法，强行改变经济规律，盲目追求非常高的生产目标，例如，全国老百姓一起制造钢铁。国家也发动了“人民公社化运动”，这个运动不符合当时农村的实际生产能力，强行把所有的财产变成集体所有，强行让所有人得到一样的收入，想用这种方法快速实现共产主义。这些做法不符合经济发展的基本道理，导致国家经济遇到巨大的困难。后来国家又发生了“文化大革命”，当时领导人错误地判断了国内的主要问题，把阶级斗争看得太重要，阶级斗争影响了所有的事情。这给党、国家和全国人民带来了巨大的灾难，中国寻找社会主义道路遇到了建国后最大的困难\upcite{ref9}。

\subsection{曲折中蕴含的“自主性”成就与制度韧性}
但是，我们用历史的眼光看待这段时间，这二十年并不是什么成绩都没有。在这个时期，美国和苏联两个大国严格封锁中国，中国承受了巨大的压力。一九六零年，苏联单方面取消了所有的合作文件，叫回了所有的技术人员，中国没有办法，只能完全靠自己发展，主动地把“自己做主，依靠自己”的原则做到最好。

中国在非常困难的情况下，依靠自己的力量取得了一系列重大成就\upcite{ref10}：

\begin{enumerate}[leftmargin=2em, itemsep=0.3em, label=\textbf{\arabic*.}]
  \item \textbf{建立独立完整的工业体系}——依靠自身力量建设国家，建立了独立的工业体系和比较完整的国民经济体系。
  \item \textbf{“三线建设”优化工业布局}——为应对可能发生的战争，在中西部地区大规模建设工厂，客观上改变了内陆工业薄弱的状况。
  \item \textbf{“两弹一星”突破大国核垄断}——在完全没有外国技术帮助的情况下，成功制造原子弹、导弹和人造卫星，国防科技取得重大突破，打破了大国对核武器的垄断，证明了中国自主建设的能力。
\end{enumerate}

\begin{summarybox}
虽经历”大跃进”和”文化大革命”等严重挫折，但中国在逆境中依靠自身力量，建立了独立工业体系，完成了”两弹一星”等重大工程，为后来的改革开放奠定了物质基础和精神财富。
\end{summarybox}

\section{新时期的成功转化：中国特色社会主义道路的确立}
中国经历了二十多年的探索，积累了很多经验，中国共产党充分认识了建设社会主义的规律。二十世纪七十年代末，面临巨大的危险，决定实行改革开放，完全放弃了传统的苏联做法，成功建立了中国特色社会主义道路。

\subsection{思想路线的拨乱反正与伟大觉醒}
重大的行动改变总是从改变想法开始，一九七八年，中国开展了一场大讨论，人们讨论如何检验真理。这场讨论批评了“文化大革命”以来的错误想法，长期以来，死板的教条限制了人们的想法，人们盲目相信苏联的做法，这场讨论极大地改变了人们的想法\upcite{ref11}。

在那一年的年底，中共召开了十一届三中全会，会议坚决停止使用“把阶级斗争放在第一位”的口号，决定把党和国家的工作重点放到经济建设上，决定实行改革开放，这是一个重大的历史决定。会议重新确定了指导思想，这个指导思想是“解放思想，根据实际情况做事”，中国确立了这个指导思想。这说明中国共产党在指导思想上重新找回了自己做主的能力，不再依靠任何外国的做法，找到一条全新的道路，为中国寻找新道路提供了坚实的认识基础。

\subsection{突破传统计划经济的藩篱：从体制改革到市场经济}
中国建立中国特色社会主义道路，在经济上，放弃了苏联式的计划经济，建立了社会主义市场经济。这个改变从农村开始，安徽省小岗村进行了尝试，农民实行了”家庭联产承包责任制”\footnote{1978年冬，安徽省凤阳县小岗村18户农民秘密签订”大包干”契约，率先实行分田到户、自负盈亏，成为中国农村改革的标志性事件。这一做法随后在全国推广，正式定名为”家庭联产承包责任制”。}，土地依然属于集体，但是农民可以自己决定如何种地。这个做法改变了过去大家收入一样的情况，农民干活更加积极，农村生产了更多的粮食\upcite{ref12}。在理论方面，中国有了新的认识，当时很多人争论计划和市场哪个属于资本主义。一九九二年，邓小平在南方视察时发表讲话，明确说明，计划只是管理经济的方法，市场也只是管理经济的方法，它们都不代表社会制度。最后，中国确定了新的经济制度，中共十四大正式提出建立“社会主义市场经济体制”，这说明中国完全放弃了苏联的做法，自己找到了管理经济的新方法\upcite{ref13}。

\subsection{“走自己的路”的理论结晶：邓小平理论与马克思主义中国化}
中国开展了改革开放的工作，邓小平带领中国共产党人总结了过去的经验，创立了邓小平理论。一九八二年，中共召开第十二次全国代表大会，邓小平在开幕式上讲话。明确指出\upcite{ref14}：

\begin{quotebox}[邓小平在中共十二大开幕式上的讲话（1982年）]
\kaishu “我们要把马克思主义的普遍真理和中国的具体实际结合起来，我们要走自己的道路，我们要建设有中国特色的社会主义，这是我们总结长期历史经验得出的基本结论。”
\end{quotebox}

邓小平说明，中国没有继续走过去死板的老路，也没有放弃社会主义制度。中国把科学社会主义的基本原则和当时的情况结合起来，结合了时代的特点，当时处于社会主义初级阶段。中国过去学习苏联，后来自己寻找道路，这个改变的过程就是马克思主义结合中国实际的第二次大发展，这个做法让中国的社会主义展现出强大的生命力。

\section{经验总结与启示：自主性在当代中国的延续与升华}
我们回顾新中国建立以来的历史，中国在建国初期必须学习苏联，后来经历困难并决定自己寻找道路。这个探索过程展现了国家发展的历史，也展现了马克思主义结合中国实际的历史，这段历史指导我们在新时代继续前进，给我们留下了宝贵的经验。

\subsection{历史逻辑：道路探索必须坚持实事求是，绝不能照搬外国模式}
历史唯物主义说明一个道理，社会生产物质的方法决定了人们的社会生活，决定了国家的政治生活，决定了人们的精神生活。苏联的做法在特定的时间里有它的好处，帮助苏联取得了一些成绩，但是，苏联的做法有一个严重的问题。苏联把一种只适合当时当地的管理制度看作永远正确的制度，把它变成了不能改变的教条。二十世纪末，苏联解体了，东欧国家也发生了巨大的变化，这些悲惨的事件证明了一个道理，一个国家不能把发展的希望放在照抄其他国家的做法上。国家照抄苏联的计划经济不能成功，照抄西方的自由主义经济也不能成功，最终一定会失败。中国的发展道路越来越宽广，它的根本原因是中国共产党一直根据实际情况做事。中国共产党根据国家真实的生产能力做事，中国共产党不断改变死板的旧制度，中国共产党自己改变自己，中国共产党自己寻找新的发展道路\upcite{ref15}。

\vspace{0.5em}
\begin{table}[htbp]
\centering
\caption{苏联模式与中国特色社会主义道路的对比}
\zihao{5}
\begin{tabularx}{0.92\textwidth}{>{\bfseries\centering\arraybackslash}p{2cm} X X}
\toprule
\textbf{对比维度} & \textbf{苏联模式} & \textbf{中国特色社会主义} \\
\midrule
经济体制 & 高度集中的计划经济，排斥市场 & 社会主义市场经济，市场起决定性作用 \\
\addlinespace
产业政策 & 片面优先发展重工业，牺牲农业 & 农业、轻工业、重工业协调发展 \\
\addlinespace
管理方式 & 中央集权，死板命令管理一切 & 央地合理分权，灵活调整政策 \\
\addlinespace
发展理念 & 照搬教条，僵化不变 & 实事求是，因地制宜，不断创新 \\
\addlinespace
历史结果 & 苏联解体，东欧剧变 & 国家持续发展繁荣，走向民族复兴 \\
\bottomrule
\end{tabularx}
\end{table}

\subsection{理论逻辑：“自主性”是马克思主义同中国实际相结合的灵魂}
我们从理论角度分析，马克思主义绝对不是不能改变的规定，它是一个开放的知识系统，随着人们的实际工作不断发展。“自己做主”是毛泽东思想的重要内容，也是邓小平理论的主要思想，也是中国特色社会主义理论的主要思想。国家提出“走自己的路”，这就要求国家遵守科学社会主义的基本原则，同时，增加中国民族的特点，增加当代时代的特点。这种自己做主的做法说明党对理论有深刻的认识，党要求我们在处理复杂问题时，不盲目相信书本知识，党要求我们不盲目听从上级命令，党要求我们只看重实际情况，党把解决中国实际问题作为发展理论的原因和目的。中国共产党一直坚持自己发展理论，所以能够冷静处理各种危险和困难，坚持自己的发展目标。

\subsection{时代启示：以“中国式现代化”全面推进中华民族伟大复兴}
时间推进到新时代，中国处在新的历史阶段，自己寻找道路的做法取得了巨大的成绩，发展出了丰富的理论，做出了很多实际工作，这个成绩就是中国式现代化\footnote{"中国式现代化"作为一个完整的理论概念，由习近平总书记在中国共产党第二十次全国代表大会报告（2022年10月）中系统阐述，明确了中国式现代化的中国特色、本质要求和重大原则。}。

中国式现代化有五个明显的特点，它们彻底改变了人们"现代化就是西方化"的错误认识\upcite{ref16}。

\vspace{0.5em}
\begin{table}[htbp]
\centering
\caption{中国式现代化的五大特征}
\zihao{5}
\begin{tabularx}{0.92\textwidth}{>{\bfseries\centering\arraybackslash}p{3.8cm} X}
\toprule
\textbf{特征} & \textbf{内涵说明} \\
\midrule
人口规模巨大的现代化 & 超大规模人口共同参与现代化建设，发展任务艰巨而独特 \\
\addlinespace
全体人民共同富裕的现代化 & 不让任何人掉队，持续缩小贫富差距，实现共同繁荣 \\
\addlinespace
物质文明和精神文明相协调的现代化 & 物质生活与精神文化同步发展，追求全面进步 \\
\addlinespace
人与自然和谐共生的现代化 & 坚持绿色发展理念，保护生态环境，实现可持续发展 \\
\addlinespace
走和平发展道路的现代化 & 不搞战争掠夺，坚持合作共赢，为世界和平作贡献 \\
\bottomrule
\end{tabularx}
\end{table}西方国家过去发动战争，抢夺其他国家的财富，用这些坏方法实现现代化，中国绝不走西方国家的旧路。苏联的做法不关心老百姓的生活，破坏了自然环境，中国绝不犯苏联犯过的错误。

中国过去学习苏联，中国后来自己寻找道路，中国今天建设中国式现代化，中国创造了新的人类文明。中国共产党的历史向全世界说明一个道理，世界上每一个国家都有权利自己寻找道路，世界上每一个民族都有能力自己寻找道路。国家和民族必须根据自己的情况做事，国家和民族必须自己寻找适合自己的现代化做法，国家必须相信自己的道路，国家必须依靠自己的力量发展自己，这就是中国建设社会主义的历史留给我们的最重要经验。

% --- 致谢部分 ---
\vspace{1.5em}
\section*{致\quad 谢}
\addcontentsline{toc}{section}{致\quad 谢}
{\kaishu
本文的完成得益于社会主义发展史课程授课教师的悉心指导与启发，老师在课堂上对社会主义发展历程的系统讲授，为本文的选题和写作提供了重要的学术方向。同时，感谢与同学们在课堂讨论和课后交流中碰撞出的思想火花，使本文的论证更加充实。在论文撰写过程中，笔者参阅了大量前辈学者的研究成果，在此一并致以诚挚的谢意。
}

\vspace{1em}
% --- 参考文献部分 ---

\begingroup
\newpage
\renewcommand{\refname}{参考文献} 
\zihao{5} 
\linespread{1.2}\selectfont 
\begin{thebibliography}{99}
\setlength{\itemsep}{0.3em} 
\setlength{\parsep}{0pt}    

\bibitem{ref1}
刘建武. 从“以苏为师”到“走自己的路”——中国共产党探索社会主义建设道路的历史演进[J]. 社会主义研究, 2021(04): 1-8.

\bibitem{ref2}
张树华. 斯大林模式的政治经济学批判及其历史教训[J]. 世界社会主义研究, 2018(06): 22-30.

\bibitem{ref3}
李瑞. 20世纪中叶国际格局演变与中国社会主义工业化起步[J]. 世界社会主义研究, 2020(03): 45-52.

\bibitem{ref4}
林毅夫, 蔡昉, 李周. 中国的奇迹: 发展战略与经济改革[M]. 上海: 上海三联书店, 上海人民出版社, 1999.

\bibitem{ref5}
逄先知, 金冲及. 毛泽东传(1949—1976)[M]. 北京: 中央文献出版社, 2003.

\bibitem{ref6}
毛泽东. 毛泽东文集: 第7卷[G]. 北京: 人民出版社, 1999.

\bibitem{ref7}
辛向阳. 科学社会主义原则在《论十大关系》中的中国化表达[J]. 科学社会主义, 2020(01): 12-19.

\bibitem{ref8}
胡绳. 中国共产党的七十年[M]. 北京: 中共党史出版社, 1991.

\bibitem{ref9}
中共中央文献研究室. 关于建国以来党的若干历史问题的决议[M]. 北京: 人民出版社, 1981.

\bibitem{ref10}
中共中央党史和文献研究院. 中国共产党简史[M]. 北京: 人民出版社, 中共党史出版社, 2021.

\bibitem{ref11}
郑异凡. 真理标准大讨论与中国特色社会主义道路的开启[J]. 当代世界与社会主义, 2018(02): 34-41.

\bibitem{ref12}
陈学明. 改革开放初期中国共产党突破传统体制藩篱的理论与实践[J]. 当代世界与社会主义, 2019(05): 55-63.

\bibitem{ref13}
萧冬连. 探路之役: 1978—1992年的中国经济改革[M]. 北京: 社会科学文献出版社, 2019.

\bibitem{ref14}
邓小平. 邓小平文选: 第3卷[G]. 北京: 人民出版社, 1993.

\bibitem{ref15}
王伟. 试论中国特色社会主义道路探索的自主性逻辑[J]. 社会主义研究, 2022(02): 15-21.

\bibitem{ref16}
姜辉. 中国式现代化蕴含的科学社会主义真理力量[J]. 科学社会主义, 2023(01): 5-12.

\end{thebibliography}
\endgroup

\end{document}